% paper/sections/07b_reinitialization.tex
% §6 続き：CLS 再初期化（Ridge-Eikonal 主軸；Split 法は legacy paragraph）・安定性制約
%
% ─────────────────────────────────────────────────────────────────────────────
\subsection{CLS 再初期化：Ridge-Eikonal 主軸（ch13 production stack）}
\label{sec:cls_compression}
% ─────────────────────────────────────────────────────────────────────────────

\noindent\textbf{Part 1 からの継承：}
本節は §~\ref{sec:reinit}（§3.2 末，再初期化方程式 \eqref{eq:cls_reinit_iso}）の
連続定式化と §~\ref{sec:ridge_eikonal}（§3.4 Ridge-Eikonal 連続定式化）に基づく
離散実装を示す．ch13 production stack の主軸再初期化は Ridge-Eikonal 法であり，
Eikonal PDE による距離関数構築 + Ridge ガウス射影で $\psi$ を再構成する．
本節の構成は以下である：
\begin{itemize}
  \item \textbf{主軸 Ridge-Eikonal 法}：§~\ref{sec:xi_sdf_reinit}（後述 $\xi$ 空間符号距離関数法）
        が ch13 production の主実装．連続定式化の理論的根拠は §~\ref{sec:ridge_eikonal} 参照．
  \item \textbf{過渡 Eikonal 法}：§~\ref{sec:eikonal_reinit}（後述「統一再初期化：Eikonal 法」）は
        Ridge-Eikonal の前段である Eikonal 距離構築の単独運用形．
\end{itemize}
旧実装（圧縮--拡散演算子分裂による再初期化）の歴史的経緯および詳細な
散逸演算子設計は §~\ref{sec:dissipative_ccd}（§4.4）に集約する．
本節以降は ch13 production である Ridge-Eikonal を主軸とする．

% 旧 SplitReinitializer paragraph (~80 LOC) は CHK-222 で削除．
% 比較参照ラベルは外部参照保護のため phantomsection alias で保持．
\phantomsection\label{sec:cls_compression_legacy}% ch13 production 非採用; ref 保護
\phantomsection\label{alg:cls_reinit_dccd}% §10_full L312 ref 保護


% ─────────────────────────────────────────────────────────────────────────────
\subsubsection{再初期化演算子の質量保存性}
\label{sec:dccd_mass_conservation}% backward-compat: 11a/11c/11d/07_advection ref 保護
% ─────────────────────────────────────────────────────────────────────────────

CLS 移流方程式（式~\eqref{eq:cls_advection}）の保存形
$\partial\psi/\partial t + \bnabla\cdot(\psi\bm{u}) = 0$ は，
領域積分により $\mathrm{d}/\mathrm{d}t \int\psi\,\mathrm{d}V = 0$
（周期／壁面 BC）を保証する．離散化においても，
空間演算子が $\sum_i [\bnabla\cdot(\psi\bm{u})]_i = 0$ を満たせば
質量は厳密に保存される．

\textbf{CCD 周期系の零和性．}
周期境界条件下のブロック巡回行列 CCD 系では，
RHS の Eq-I（1階微分）が反対称ステンシル $\propto (f_{i+1} - f_{i-1})$，
Eq-II（2階微分）が対称ステンシル $\propto (f_{i-1} - 2f_i + f_{i+1})$ を持ち，
ともに全格子点の和が厳密に零となる．
連立方程式 $A\bm{x} = \bm{b}$ の $\bm{b}$ が零和であり $A$ が正則であるから，
$\sum_i d_{1,i} = 0$，$\sum_i d_{2,i} = 0$ が厳密に成立する%
\footnote{RHS の各成分は反対称/対称ステンシルから構成され，
周期和は零．非特異な係数行列を通しても零和性は保存される．}．
帯状散逸付き 1 次微分演算子も 2 階差分のテレスコーピング和が
零となるため，全点和に対して同様の零和性が保たれる．

\textbf{Ridge-Eikonal 主軸での運用．}
Ridge-Eikonal 法は反復を伴わない 1 ショット再構成
（§~\ref{sec:xi_sdf_reinit}）であり，演算子分裂による分布的質量損失は
構造的に発生しない．離散総量の補正は事後的な体積重み付き
$\phi$ 空間質量補正（§~\ref{sec:dgr_thickness} の補正手順）で行う．
旧実装の演算子分裂解法における質量損失メカニズムと
散逸演算子レベルの保存性議論の詳細は，§~\ref{sec:dissipative_ccd}
（§4.4 散逸 CCD の歴史的経緯）に集約する．

\textbf{非一様格子への拡張．}
非一様格子ではセル体積 $\Delta V_{i,j} = h_x(i) \cdot h_y(j)$ が
格子点ごとに異なるため，質量積分を
$M = \sum_{i,j} \psi_{i,j}\,\Delta V_{i,j}$ に置き換え，
界面重み $W = \sum w_{i,j}\,\Delta V_{i,j}$ も体積重み付きとする．
一様格子では $\Delta V$ が定数に退化するため，従来の無重み和に一致する．
非一様格子では計算空間零和性が物理空間の体積重み付き和へ
自動的に移らないため，clamp・再初期化・格子転写後に
体積重み付き質量補正を適用して総量を所定の許容誤差へ戻す運用とする．


% ─────────────────────────────────────────────────────────────────────────────
\subsubsection{適応的再初期化トリガー}
\label{sec:adaptive_reinit}
% ─────────────────────────────────────────────────────────────────────────────

固定頻度の再初期化（例：10 ステップ毎）は過剰な再初期化を引き起こし，
形状精度を大幅に劣化させることが数値実験で判明した
（§~\ref{sec:verify_dccd_conservation}）．
再初期化スキームに帯状散逸を併用する場合の高周波減衰は
界面プロファイルの主要波長に対して僅かであり，形状誤差への寄与は小さい．
支配的な形状誤差源は，各再初期化呼び出しにおける
界面再構成ステップの累積的プロファイル歪みである．

そこで，再初期化の呼び出しを体積モニタ $M(\tau) = \int \psi(1-\psi)\,\mathrm{d}V$
に基づく適応的トリガーに置き換える：
\begin{equation}
  \text{再初期化の実行条件：}\quad
  \frac{M(\tau)}{M_{\mathrm{ref}}} > \theta_{\mathrm{reinit}}
  \label{eq:adaptive_reinit_trigger}
\end{equation}
ここで $M_{\mathrm{ref}}$ は直前の再初期化直後の $M$ 値，
$\theta_{\mathrm{reinit}}$ は界面太り度の許容比を与える閾値パラメタ
（既定値の選定および感度は §~\ref{sec:val_summary} を参照）である．
$M(\tau)$ は界面プロファイルの太り度を表し，
$\psi = H_\varepsilon(\phi)$ の平衡形からの逸脱に比例して増大する．
適応トリガー導入による再初期化回数の削減効果と形状誤差改善は
§~\ref{sec:verify_dccd_conservation} を参照．


% ─────────────────────────────────────────────────────────────────────────────
\subsubsection{演算子分裂の構造的界面膨張と Direct Geometric 厚み補正}
\label{sec:dgr_thickness}
% ─────────────────────────────────────────────────────────────────────────────

\textbf{分裂誤差による界面膨張．}
旧実装の演算子分裂解法（§~\ref{sec:dissipative_ccd}）は，
圧縮（陽解法 FE）と拡散（CN-ADI）を\textbf{逐次適用}する．
平衡では両項が釣り合うが，逐次実行では：
(1) 圧縮が先に界面を鋭化 $\varepsilon' < \varepsilon$，
(2) 拡散が鋭化後のプロファイルに対して過大に作用 $\varepsilon'' > \varepsilon$，
という $\Ord{\Delta\tau}$ の分裂誤差が生じる．
数値検証により，\textbf{完全な平衡プロファイル} $\psi = H_\varepsilon(\phi)$
（移流なし）に再初期化を適用するだけで，
1回あたり $\varepsilon_{\mathrm{eff}}/\varepsilon \approx 1.4$
の膨張が確認された（§~\ref{sec:verify_dccd_conservation}）．
89回の累積で $\varepsilon_{\mathrm{eff}} \approx 4\varepsilon$ に達する．
$n_{\mathrm{reinit}}$ を増加させると分裂反復が増え，膨張は悪化する．

\textbf{Direct Geometric Reinitialization（DGR）．}
この構造的膨張を補正するため，logit 逆変換と勾配正規化に基づく
厚み補正を形状再初期化とは\textbf{独立に}適用する：
\begin{enumerate}
  \item 実効厚の推定：界面帯域（$0.05 < \psi < 0.95$）で
    $\varepsilon_{\mathrm{eff}} = \mathrm{median}\!\left[\psi(1-\psi)/|\bnabla\psi|\right]$
  \item 逆変換：$\phi_{\mathrm{raw}} = \varepsilon\,\ln[\psi/(1-\psi)]$
  \item 再スケール：$\phi_{\mathrm{sdf}} = \phi_{\mathrm{raw}} \cdot (\varepsilon_{\mathrm{eff}}/\varepsilon)$
    （$\varepsilon_{\mathrm{eff}}>\varepsilon$ では $\phi_{\mathrm{raw}}
    = \varepsilon\ln[\psi/(1-\psi)] = (\varepsilon/\varepsilon_{\mathrm{eff}})\phi_{\mathrm{true}}$
    と縮小しているため，$\varepsilon_{\mathrm{eff}}/\varepsilon$ を掛けて真の SDF を復元する）
  \item 再構成：$\psi_{\mathrm{new}} = H_\varepsilon(\phi_{\mathrm{sdf}})$
  \item $\phi$ 空間質量補正：
    $\phi_{\mathrm{sdf}} \leftarrow \phi_{\mathrm{sdf}} + \delta\phi$，
    $\psi_{\mathrm{new}} = H_\varepsilon(\phi_{\mathrm{sdf}})$，
    ここで $\delta\phi = \delta M / \displaystyle\int H'_\varepsilon\,\mathrm{d}V$，
    $H'_\varepsilon = \psi_{\mathrm{new}}(1-\psi_{\mathrm{new}})/\varepsilon$（$= \mathrm{d}\psi/\mathrm{d}\phi$）
\end{enumerate}

この補正は，$\psi = H_{\varepsilon_{\mathrm{eff}}}(\phi_{\mathrm{true}})$ の
形のプロファイルに対して厳密に $\psi_{\mathrm{new}} = H_\varepsilon(\phi_{\mathrm{true}})$
を復元する（DGR 復元性：$\phi_{\mathrm{raw}} = (\varepsilon/\varepsilon_{\mathrm{eff}})\phi_{\mathrm{true}}$ と
$|\bnabla\phi_{\mathrm{raw}}| = \varepsilon/\varepsilon_{\mathrm{eff}}$ より）．
$\phi$ 空間質量補正は界面を一様に $\delta\phi$ だけシフトする；
DGR が $|\bnabla\phi_{\mathrm{sdf}}| \approx 1$ を保証するので
界面の物理的変位 $\delta x = \delta\phi / |\bnabla\phi| \approx \delta\phi$ は
曲率非依存の一様シフトとなり形状を変えない（DGR 幅保存性）．

\textbf{注意：} 旧実装では $\psi$ 空間補正
$\psi \leftarrow \psi_{\mathrm{new}} + (\delta M / \sum w)\,w$，
$w = 4\psi_{\mathrm{new}}(1-\psi_{\mathrm{new}})$ を使用していた．
この補正による界面変位は $\delta x \propto w/|\bnabla\psi|$ であり，
曲率が大きい領域（$|\bnabla\psi|$ 小）ほど変位が大きい．
毛管振動液滴のような高曲率界面では，
各再初期化呼び出しで系統的に mode-2 変形が増幅し $D(t)$ が正しく減衰しない．
$\phi$ 空間補正はこの質量保存の問題を解決する（体積保存誤差：$347\% \to 0.02\%$）．
ただし $D(t)$ の誤差は解消しない――根本原因は質量補正ではなくシャープニングの非一様性である
（§~\ref{sec:val_capillary}参照）．

\textbf{運用：}
DGR は形状修復（圧縮--拡散）とは独立に，物理時間 20 ステップ毎に適用する．
形状再初期化が界面プロファイルを tanh 形に整形し，
DGR がその直後（または独立に）界面厚を $\varepsilon$ に復元する．
単一渦テスト（$N=128$, $\varepsilon/h=1.0$）で
$\varepsilon_{\mathrm{eff}}/\varepsilon \approx 1.03$ を全時刻で維持し，
面積誤差を従来比 24 倍改善した（§~\ref{sec:verify_dccd_conservation}）．

\textbf{適用制限（1）：幾何形状．}
DGR のグローバル $\varepsilon_{\mathrm{eff}}$ 推定は
界面厚の膨張と狭い形状特徴（スロット等）を区別できない．
Zalesak スロット円板のように界面幅 $\sim 6\varepsilon$ の狭い構造を含む場合，
DGR がスロット内部を液相に充填してしまうため\textbf{適用してはならない}．
鋭い幾何形状には従来の圧縮--拡散再初期化（固定頻度）を使用する．

\textbf{適用制限（2）：$\sigma > 0$ 毛管波ダイナミクス（界面折れ）．}
表面張力が存在する系では，毛管移流により界面帯域内に
$|\bnabla\psi| \to 0$ となる局所的な折れ（fold）が生じうる．
DGR の中央値 $\varepsilon_{\mathrm{eff}}$ 推定は
fold セル（局所異常値）に対してロバストであるため，
fold を検出できず事実上何もしないまま返す．
次ステップの CCD Laplacian が fold セルの上で非物理的な曲率スパイクを生成し，
CSF 力 $\sigma\kappa\delta_\varepsilon\bnabla\psi$ が指数的に増大して発散する．

数値実験（§~\ref{sec:val_capillary}，Prosperetti 1981 ベンチマーク）により：
\begin{itemize}
  \item DGR 単独：$t < 0.2$ で発散（運動エネルギー $> 10^6$）
  \item ハイブリッド（分裂再初期化 + DGR，$\phi$ 空間補正）：発散は回避，体積保存は正確，
    しかし $D(t)$ が誤った値に飽和．
    原因：DGR のグローバル中央値 $\varepsilon_{\mathrm{eff}}$ による一様スケールが
    曲率分布に対して非一様なシャープニングを引き起こし（圧縮端は過剰，伸展端は過小），
    mode-2 変形を系統的に増幅する
  \item \textbf{分裂再初期化単独：正しい $D(t)$ を与え $T=10$ まで安定}（推奨）
\end{itemize}
$\sigma > 0$ の毛管波・液滴振動ベンチマークには
\textbf{分裂再初期化のみ}を用いる（DGR の後段追加は不可）．


% ─────────────────────────────────────────────────────────────────────────────
\subsubsection{Ridge-Eikonal 再初期化：Eikonal 基盤と Ridge 補正の統合定式化}
\label{sec:xi_sdf_reinit}
\phantomsection\label{sec:eikonal_reinit}% 旧 §6.3.4 ラベル統合; 外部参照保護
% ─────────────────────────────────────────────────────────────────────────────

\noindent\textbf{Part 1 連続定式化との関係：}
本節は §~\ref{sec:reinit}（§3.3 再初期化方程式）と
§~\ref{sec:ridge_eikonal}（§3.4 Ridge-Eikonal 連続定式化）の離散実装である．
ch13 production stack の主軸再初期化はここで定式化する Ridge-Eikonal 法
（前段の Eikonal 距離構築 + Ridge ガウス射影による $\psi$ 再構成）であり，
旧実装（圧縮--拡散演算子分裂）は §~\ref{sec:dissipative_ccd}（§4.4）に集約する．

\paragraph*{Eikonal 基盤と継承される理論保証}
Ridge-Eikonal は Eikonal PDE による符号距離関数構築を前段に持つ．
本項では Ridge-Eikonal が継承する Eikonal 法の基盤と理論保証を整理する．
Eikonal 方程式 \cite{Sussman1994}
\begin{equation}
  \frac{\partial\phi}{\partial\tau} + \operatorname{sgn}(\phi_0)\bigl(|\nabla\phi| - 1\bigr) = 0
  \label{eq:eikonal_pde}
\end{equation}
は零点集合（界面位置）を不変に保ちつつ $|\nabla\phi| \to 1$ を回復する（符号距離関数化）．
これによりフォールド（$|\nabla\phi|\to 0$ 領域）も自然に修正される．

\begin{tcolorbox}[algbox, title={Eikonal 再初期化アルゴリズム}]
\begin{enumerate}
  \item \textbf{ロジット反転：}$\phi = \varepsilon\ln(\psi/(1-\psi))$（飽和クランプ付き）
  \item \textbf{Eikonal PDE 反復（疑似時間反復；反復回数は §~\ref{sec:val_summary} 参照）：}
        Godunov 一次風上差分で式~\eqref{eq:eikonal_pde} を疑似時間積分
        （\S\ref{sec:scheme_per_variable} 表で $\phi$ 再初期化に WENO--HJ を推奨する記述は
        Level 2/3 の高精度オプション；Level 1 検証では Godunov で十分）
        \begin{equation}
          \phi^{(k+1)} = \phi^{(k)} - \Delta\tau\,\operatorname{sgn}(\phi_0)\,
          \bigl(G^{\mathrm{Godunov}}(\phi^{(k)}) - 1\bigr)
        \end{equation}
  \item \textbf{セル局所 $\varepsilon$ による再構成：}
        $\psi(i,j) = H_{\varepsilon(i,j)}(\phi)$，
        $\varepsilon(i,j) = \varepsilon_\xi\max(h_x(i),\,h_y(j))$

        一様格子では $\varepsilon(i,j) = \varepsilon$（スカラーに帰着）；
        非一様格子では各セルの物理格子幅に比例した界面幅を実現する．
  \item \textbf{$\phi$ 空間質量補正：}
        $\delta\phi = \Delta M / \int H'_\varepsilon\,dV$，
        $\phi \leftarrow \phi + \delta\phi$，$\psi \leftarrow H_{\varepsilon}(\phi)$
\end{enumerate}
\end{tcolorbox}

\noindent\textbf{理論的保証：}
\begin{enumerate}
  \item \textbf{零点集合保存（正確）：}Eikonal PDE はレベル集合の輸送のみを行い，
        $\phi = 0$ 面を不変に保つ→界面位置（形状）は変化しない．
  \item \textbf{勾配収束：}十分な反復で $|\nabla\phi|\to 1$ が保証され，
        界面幅は $\varepsilon(i,j)$ に正確に収束する（DGR のグローバル中央値近似不要）．
  \item \textbf{質量保存（前クリッピング）：}Step 4 の $\phi$ 空間補正
        $\delta\phi=\Delta M/\!\int H'_\varepsilon\,dV$ により設計上
        $\sum\psi_{\mathrm{corr}}=\sum\psi_{\mathrm{old}}$ が成立する
        （$H'_\varepsilon$ の正値性と $\delta\phi$ の構成より直ちに従う；
        クリッピング前段で評価する場合に限る）．
\end{enumerate}

\paragraph*{Eikonal 単独運用の限界：Godunov 由来の零セットドリフト}
Eikonal 法は $\sigma{=}0$ の受動移流・静的形状問題（Zalesak スロット等）では
体積保存・界面幅・形状復元のすべての条件を満たす有力手法であるが，
$\sigma>0$ で界面振動を伴う系では Godunov 差分の離散誤差により
零点集合が各呼び出しでわずかに移動し，反復累積で形状誤差を生む．
これは DGR のグローバル中央値非一様性とは異なるが，同等の系統的誤差として現れる
（定量比較は §~\ref{sec:val_capillary} を参照）．
本制約は $\phi$ の局所操作を反復する全手法に共通であり，
\textbf{反復を一切伴わない解析的距離構築}を要請する——これが
Ridge-Eikonal（次項 ξ 空間符号距離関数法）の動機である．


\paragraph*{Ridge-Eikonal による解：ξ 空間符号距離関数法}
反復累積誤差を原理的に排除するため，
ch13 production stack では再初期化呼び出しごとに
現在の $\psi$ 場から零点交差を直接検出し，
格子インデックス空間（$\xi$ 空間）における符号距離を
解析的に構築する \textbf{ξ 空間符号距離関数（ξ-SDF）法}を採用する．
本手法は §~\ref{sec:ridge_eikonal}（§3.4 Ridge-Eikonal 連続定式化）の Ridge ガウス
射影を離散化したものであり，反復を行わないため累積ドリフトが構造的に存在しない．

\begin{tcolorbox}[algbox, title={ξ-SDF 再初期化アルゴリズム (Strategy B)}]
\begin{enumerate}
  \item \textbf{ロジット反転：}
        $\phi_{\mathrm{raw}}(i,j) = \varepsilon\ln\!\bigl(\psi(i,j)/(1{-}\psi(i,j))\bigr)$
        （飽和クランプ付き）
  \item \textbf{零点交差の検出（線形補間）：}$\phi_{\mathrm{raw}}$ の符号が変化する隣接セル対を探索し，
        格子座標上の交差位置を線形補間で決定する．
        $x$-方向（$(i,j)\to(i{+}1,j)$）の交差 $\xi$-座標：
        \begin{equation}
          \xi^* = i + \frac{|\phi_{\mathrm{raw}}(i,j)|}{|\phi_{\mathrm{raw}}(i,j)| + |\phi_{\mathrm{raw}}(i{+}1,j)|}
          \label{eq:xi_crossing_x}
        \end{equation}
        $y$-方向（$(i,j)\to(i,j{+}1)$）の交差 $\eta$-座標：
        \begin{equation}
          \eta^* = j + \frac{|\phi_{\mathrm{raw}}(i,j)|}{|\phi_{\mathrm{raw}}(i,j)| + |\phi_{\mathrm{raw}}(i,j{+}1)|}
          \label{eq:xi_crossing_y}
        \end{equation}
        全交差位置の集合を $S_\xi = \{(\xi^*_k,\,\eta^*_k)\}_{k=1}^{N_c}$ とする．
  \item \textbf{ξ距離場の構築：}各セル $(i,j)$ から最近傍零点交差までの
        ξ空間 Euclid 距離を計算し，$\phi_{\mathrm{raw}}$ の符号を付与する：
        \begin{equation}
          \phi_\xi(i,j) = \operatorname{sgn}\!\bigl(\phi_{\mathrm{raw}}(i,j)\bigr)
          \cdot \min_{(\xi^*,\eta^*)\in S_\xi}
          \sqrt{(i-\xi^*)^2 + (j-\eta^*)^2}
          \label{eq:xi_sdf}
        \end{equation}
  \item \textbf{$\psi$ 再構成：}定数 $\varepsilon_\xi = \varepsilon/h_{\min}$
        （インターフェース幅を格子幅比で表したスカラー定数）を用いて
        \begin{equation}
          \psi_{\mathrm{new}}(i,j) = H_{\varepsilon_\xi}\!\bigl(\phi_\xi(i,j)\bigr)
          = \frac{1}{1+\exp\!\bigl(-\phi_\xi(i,j)/\varepsilon_\xi\bigr)}
          \label{eq:xi_sdf_psi}
        \end{equation}
  \item \textbf{$\phi_\xi$-空間質量補正：}
        $\delta\phi = \Delta M / \int H'_{\varepsilon_\xi}\,dV$，
        $\phi_\xi \leftarrow \phi_\xi + \delta\phi$，
        $\psi \leftarrow H_{\varepsilon_\xi}(\phi_\xi)$
\end{enumerate}
\end{tcolorbox}

\paragraph*{理論的性質}

\begin{description}
\item[命題~1（零点集合の完全保存）]
$\phi_{\mathrm{raw}}(i,j) = 0
  \;\Longleftrightarrow\; \phi_\xi(i,j) = 0
  \;\Longleftrightarrow\; \psi_{\mathrm{new}}(i,j) = 0.5$

\textbf{前提}：本命題は\textbf{$\mathrm{sgn}(0)\equiv 0$}の取決めに従う；
$\phi_{\mathrm{raw}}(i,j)=0$ のとき隣接交差検出
（式~\eqref{eq:xi_crossing_x}）は分子 $|\phi_{\mathrm{raw}}(i,j)|$ が
$0$ となる退化を許容し，自身を $S_\xi$ の元として加える．

\textit{証明：}$\phi_{\mathrm{raw}}(i,j)=0$ のとき，
式~\eqref{eq:xi_crossing_x} の補間分子は 0 となり交差位置 $\xi^* = i$ が得られる
（すなわち点 $(i,j)$ 自身が $S_\xi$ の元となる）．
よって式~\eqref{eq:xi_sdf} の最小値は 0 となり $\phi_\xi(i,j) = 0$，
式~\eqref{eq:xi_sdf_psi} より $\psi = 1/(1+e^0) = 0.5$．逆方向も定義より自明．\qed

\item[命題~2（ξ-SDF 性質：$|\nabla_\xi\phi_\xi|=1$）]
式~\eqref{eq:xi_sdf} で定義される $\phi_\xi$ は，格子インデックス空間において
$|\nabla_\xi\phi_\xi|=1$ を満たす真の符号距離関数である（Voronoi 境界を除く）．

\textit{証明：}$d(i,j;S_\xi) = \min_k\sqrt{(i-\xi^*_k)^2+(j-\eta^*_k)^2}$ は
$\xi$ 空間の点集合 $S_\xi$ への Euclid 距離場であり，
$|\nabla d|=1$ はほぼすべての点で成立する（距離場の基本性質）．
$\phi_\xi = \operatorname{sgn}(\phi_{\mathrm{raw}}) \cdot d$ であり，$\operatorname{sgn}$ の符号変化は
$S_\xi$ 上（$d=0$）でのみ起きるため，$|\nabla_\xi\phi_\xi|=|\nabla d|=1$．\qed

\item[命題~3（無反復 $\Rightarrow$ ドリフト累積なし）]
ξ-SDF は反復疑似時間積分を行わない．
各呼び出しの零点集合移動量は線形補間誤差 $O(h^2)$ のみであり，
Godunov 差分の $O(\Delta\tau\cdot n_{\mathrm{iter}})$ 累積ドリフトを生じない．
\end{description}

\noindent\textbf{検証結果（Prosperetti ベンチマーク，$\alpha=1.0$，$64\times64$）：}
\begin{center}
\small
\begin{tabular}{lccc}
\hline
手法 & $D(T{=}2)$ & $D(T{=}10)$ & 判定 \\
\hline
Eikonal（ZSP なし） & 0.245 & — & $\times$ \\
Eikonal + ZSP（Strategy A）   & 0.129 & — & $\times$ \\
\textbf{ξ-SDF（Strategy B，本節）} & \textbf{0.050} & \textbf{0.226} & $\times$（$T{=}10$ 未達）\\
分裂法のみ（参照）   & 0.037 & 0.0036 & $\checkmark$ \\
\hline
\end{tabular}
\end{center}

\noindent$T{=}2$ では $D=0.050$（閾値 $0.05$ 境界値），$T{=}10$ では $D=0.226$，
体積誤差 $\Delta V/V_{\max} = 3.5\%$（目標 $<1\%$ 未達）．

\noindent\textbf{Fast Marching Method による検証：}
前節 ξ-SDF の失敗原因として当初 Voronoi kink 仮説（ξ-SDF の $C^0$ 不連続が
6次 CCD で増幅 $\to$ $\kappa$ ノイズ $\to$ 体積誤差）を立てた．
これを検証するため，\cite{Sethian1996} の Fast Marching Method (FMM) を実装した
（二次 Godunov 更新 $d = \tfrac{1}{2}(a_x+a_y+\sqrt{2-(a_x-a_y)^2})$；反復なし，
単一前進波で $|\nabla_\xi\phi|=1$ を達成する $C^1$ 場を構築）．
FMM は 2次微分ノイズを低減する（標準偏差 $3.93\to 2.83$）が，
$T{=}1$ で体積誤差 $8.2\%$（ξ-SDF の 5 倍超）を示し仮説を反証した．
静的テスト（advection なし，200 回再初期化）では両手法とも体積誤差 $\approx 0\%$ であり，
誤差は advection 段階に起因することが確認された．

\noindent\textbf{修正仮説—界面幅効果：}
分裂法は拡散 PDE により界面幅が $\varepsilon_{\mathrm{eff}}\approx 1.4\varepsilon$ に広がる．
この「意図せざる広がり」が表面張力の空間集中を緩和し，
圧力ポアソン方程式の投影精度不足（$\nabla\cdot\mathbf{u}\neq 0$ 残差）による
体積誤差を抑制している（式~\eqref{eq:dt_sigma} の毛管波安定条件参照）．
ξ-SDF および FMM は正確な幅 $\varepsilon$ の界面を生成するが，
これが $\sigma>0$ 問題での数値的不安定性の原因となる．

\begin{equation}
  \Delta V / V_0 \approx \frac{\Delta t}{\rho} \int \psi\,\nabla\cdot\mathbf{u}^*\,dV
  \label{eq:volconsrate}
\end{equation}

$\nabla\cdot\mathbf{u}^*$ は未投影速度の残留発散であり，
表面張力が $\varepsilon^{-1}$ スケールで集中するほど PPE の右辺が大きくなる．
$\sigma>0$ 毛管波ベンチマークには\textbf{分裂法単独}を推奨する方針は不変であり，
ξ-SDF および FMM の適用は $\sigma=0$ 問題に限定する
（各手法の定量比較は表~\ref{tab:reinit_comparison_chk138}）．

\begin{table}[H]
  \centering
  \small
  \begin{tabular}{lcccccc}
    \hline
    方法 & 形状 & 幅 $\varepsilon$ & $\sigma>0$ & $D(T{=}2)$ & VolCons & 反復 \\
    \hline
    分裂法のみ & $\checkmark$ & ${\sim}1.4\varepsilon$ & $\checkmark$ & 0.037 & ${<}1\%$ @$T{=}10$ & あり \\
    DGR のみ & $\times$ & $\checkmark$ & $\times$ Blowup & — & — & なし \\
    ハイブリッド & $\checkmark$ & $\checkmark$ & $\times$ & 0.129 & 低 & あり \\
    Eikonal & $\checkmark$ & $\checkmark$ & $\times$ & 0.245 & 低 & あり \\
    Eikonal+ZSP（Strategy A）& $\checkmark$ & $\checkmark$ & $\times$ & 0.129 & 低 & あり \\
    ξ-SDF（Strategy B）& $\checkmark$ & $\checkmark$ & $\times$ & 0.050 & $3.5\%$ @$T{=}10$ & なし \\
    \textbf{FMM}& $\checkmark$ & $\checkmark$ & $\times$ & — & $\mathbf{8.2\%}$ @$T{=}1$ & なし \\
    \textbf{ξ-SDF（$f{=}1.4$）}& $\checkmark$ & $\checkmark$ & $\checkmark$ & \textbf{0.028} & $1.38\%$ @$T{=}2$ & なし \\
    \hline
  \end{tabular}
  \caption{再初期化手法の比較．ε幅拡大（$f{=}1.4$）が最良の $D(T{=}2)$を実現．}
  \label{tab:reinit_comparison_chk138}
\end{table}

\noindent\textbf{ε幅拡大による修正：}
前項の界面幅仮説に基づき，ξ-SDF 再構成のパラメータに $f{=}1.4$ を設定して
$\varepsilon_{\text{eff}} = 1.4\,\varepsilon$ とすることで，split-only の自然な幅広化と等価な
正則化を明示的に与える．

再初期化の Step~3 を
\begin{equation}
  \psi_{\text{new}}(i,j) = H_{f\varepsilon_\xi}(\phi_\xi(i,j)), \quad f = 1.4
  \label{eq:xi_sdf_wide}
\end{equation}
と修正する（コード実装：\texttt{EikonalReinitializer(eps\_scale=1.4)}）．
Step~4 の質量補正の重み $W = \sum H'_{f\varepsilon_\xi}$ も自動的に調整される
（$W \propto 1/(f\varepsilon_\xi)$ となるため $\delta\phi$ が増大するが，単位体積あたりの質量誤差は同等）．

\textbf{検証結果：}
\begin{center}
\begin{tabular}{lccc}
  \hline
  & $D(T{=}1)$ & $D(T{=}2)$ & VolCons max \\
  \hline
  ξ-SDF（$f{=}1.0$） & — & 0.050 & $3.5\%$ @$T{=}10$ \\
  \textbf{ξ-SDF（$f{=}1.4$）} & \textbf{0.018 $\checkmark$} & \textbf{0.028 $\checkmark$} & \textbf{1.38\% @$T{=}2$ $\triangle$}$^{\dagger}$ \\
  split-only（参考） & — & 0.037 & $<1\%$ @$T{=}10$ \\
  \hline
\end{tabular}
\end{center}

$T{=}1$ では VolCons$=0.80\%<1\%$ を達成し，$D(T{=}1)=0.018$ は split-only の
$D(T{=}2)=0.037$ より小さい．$T{=}2$ では VolCons$=1.38\%$ と目標 $1\%$ を僅かに超過するが，
ξ-SDF($f{=}1.0$) の 3.5\% に比べ約 2.5 倍の改善であり，$D(T{=}2)=0.028$ は全手法中最良値を示す．

\noindent$^{\dagger}$ 目標 VolCons $<1\%$ に対し $T{=}2$ で $1.38\%$ は部分達成（partial）の意で $\triangle$ 表記．

\textbf{残課題：}$T{=}10$ での長時間挙動は未検証．VolCons 残差の起源は，
毛管波のダイナミクスと連動した界面変形による PPE 残差の位相依存変動と考えられる
（$t{=}0.5$ → $t{=}1.5$ にかけて VolCons が一度減少する非単調挙動がこれを示唆する）．

% ─────────────────────────────────────────────────────────────────────────────
\subsubsection{再初期化の実装ガイドと収束判定}
% ─────────────────────────────────────────────────────────────────────────────

\begin{tcolorbox}[mybox, title={実装ガイド：再初期化の定量的運用}]
\textbf{仮想時間刻み：} $\Delta s = \min(h_x, h_y)$ として，
安定条件（\ref{warn:cls_dtau_stability}）の 2 次元拡散安定限界に対し
保守的な安全マージンを取った推奨値を用いる．

\medskip
\noindent\textbf{収束モニター：} $M(\tau) = \int_\Omega \psi(1-\psi)\,dV$ を使用する．
基準値
\begin{equation}
  M_{\mathrm{ref}} = \int_\Omega H_\varepsilon(\phi_0)\bigl(1-H_\varepsilon(\phi_0)\bigr)\,dV
\end{equation}
に対して $M(\tau)/M_{\mathrm{ref}}$ が閾値 $\theta_{\mathrm{reinit}}$ を超えたときのみ再初期化を実行する
（適応的再初期化；§~\ref{sec:adaptive_reinit}．既定値の選定および感度は §~\ref{sec:val_summary}）．
軽微な乱れでは3～5ステップ，界面トポロジー変化直後は10～20ステップ程度が目安．

\medskip
\noindent\textbf{終了判定（再初期化ループの停止条件）：}
再初期化ループを反復ごとに以下の条件でチェックし，いずれかを満たしたとき停止する：
\begin{enumerate}
  \item \textbf{残差収束：}
        $\|\psi^{m+1} - \psi^m\|_{L^\infty} < \delta_{\tau}$（推奨：$\delta_\tau = 10^{-4}\,\Delta s$）
  \item \textbf{ステップ数上限：}
        $m \ge N_{\tau,\max}$（推奨：$N_{\tau,\max} = 5$；CLS の lift-and-redistribute が
        1 回で $M(\tau)/M_{\mathrm{ref}}$ を閾値以下まで戻すため通常運転ではこれで十分．
        トポロジー変化直後のみ $N_{\tau,\max} = 20$ に引き上げる．
        校正値は §~\ref{sec:verification} ベンチマークに基づく）
\end{enumerate}
物理タイムステップあたりの再初期化ステップ数（1-5 程度）が少なすぎる場合は
$M(\tau)$ の時系列を記録し，蓄積的な乱れを検知すること．

\noindent\textbf{トレードオフ：}
再初期化不足 → プロファイル崩れ・曲率誤差増大；
再初期化過剰 → 界面位置の微小移動・体積誤差蓄積．
\end{tcolorbox}

% ─────────────────────────────────────────────────────────────────────────────
\subsubsection{仮想時間刻み \texorpdfstring{$\Delta\tau$}{Δτ} の安定性}
\label{warn:cls_dtau_stability}
% ─────────────────────────────────────────────────────────────────────────────
以下は，圧縮項と拡散項をともに陽的に扱う場合の参考安定条件である．
ch13 production stack は §~\ref{sec:xi_sdf_reinit} の Ridge-Eikonal を主軸とし
反復を伴わないため，本節の陽的 CFL 制約は直接の運用拘束ではなく，
旧実装（§~\ref{sec:dissipative_ccd}）との比較参照のために記す．
完全陽解法として再初期化方程式を TVD-RK3 で仮想時間積分する場合，
$\Delta\tau$ には以下の安定性制約が存在する．
再初期化方程式は双曲型（圧縮項）と放物型（拡散項）の混合型であり，
それぞれの安定条件は以下の通りである．
双曲型成分（CFL 条件）では，フラックスヤコビアン $|F'(\psi)| = |1-2\psi|$ の最大値は
$\psi \in [0,1]$（値域クランプにより保証）の両端で達成され $\max = 1$ であるから
$\Delta\tau_{\text{hyp}} \le \Delta s$（$\Delta s = \min(h_x, h_y)$）．
放物型成分（拡散安定条件）では，2次元で
$\Delta\tau_{\text{par}} \le (\Delta s)^2/(4\varepsilon)$．
実用的な安定条件はこれらの最小値：
\[
  \Delta\tau \le \min\!\left(\Delta s,\; \frac{(\Delta s)^2}{4\varepsilon}\right)
\]
典型的なパラメータ $\varepsilon \sim \Delta s$ オーダーでは
$(\Delta s)^2/(4\varepsilon) \sim \Delta s$ オーダーとなり，
拡散型成分の制約が支配的である．
\textbf{完全陽解法での参考値：}拡散安定限界に対して
保守的な安全余裕を確保した $\Delta\tau$ を採用する．
旧実装の圧縮--拡散演算子分裂解法（散逸 CCD 圧縮項 + Crank--Nicolson 拡散項）の
具体的な更新式と歴史的経緯は §~\ref{sec:dissipative_ccd}（§4.4）に集約する．
\textbf{備考（圧縮項の離散化選択）：}
CLS 再初期化圧縮項は非線形フラックス $\bnabla\cdot(\psi(1-\psi)\hat{n})$ を含む．
UCCD6（§~\ref{sec:uccd6_def}）は線形スカラー移流 $\bu\cdot\bnabla\cdot$ を想定した
特性方向組み替えを用いるため，非線形フラックスにそのまま適用すると
風上性の保証が崩れる．本稿では NS 側のスカラー移流にのみ UCCD6 を適用し，
界面再初期化は Ridge-Eikonal による反復不要の 1 ショット再構成で扱う．
UCCD6 の非線形フラックス化（Jacobian $F'(\psi) = 1-2\psi$ の符号に応じた
左右ステンシル動的切替）は本稿の範囲外であり，§~\ref{sec:future_work} で将来課題として整理する．

% ─────────────────────────────────────────────────────────────────────────────
\subsubsection{安定性制約（時間刻みの決定）}
\label{sec:stability}
% ─────────────────────────────────────────────────────────────────────────────

時間刻み $\Delta t$ は，以下の3種類の制約のうち最も厳しい条件を選ぶ（安全率 $\text{CFL}_{\max} < 1$）：

\begin{tcolorbox}[resultbox, title={時間刻みの制約（CN 法採用時：2制約）}]
\begin{align}
  &\text{\textbf{対流 CFL}：}\quad
  \Delta t_{\text{adv}} = \text{CFL}_{\max} \cdot \frac{1}{\max\!\left(\dfrac{|u|_{\max}}{\Delta x}+\dfrac{|v|_{\max}}{\Delta y}\right)}
  \label{eq:dt_adv} \\[6pt]
  &\text{\textbf{毛管波制約（2D）}：}\quad
  \Delta t_\sigma = \sqrt{\frac{(\rho_l + \rho_g)\min(\Delta x, \Delta y)^3}{2\pi\sigma}}
  \label{eq:dt_sigma}
\end{align}
\textbf{最終的な時間刻み：}$\Delta t = \min(\Delta t_{\text{adv}},\, \Delta t_\sigma)$

\medskip
\textbf{粘性 CFL（参考値）：}
\[
  \Delta t_{\text{visc}} = \frac{\min(\Delta x, \Delta y)^2}{4(\mu/\rho)_{\max}}
  \quad \text{（陽的粘性処理のみ適用；CN 法採用時は不要）}
\]
CN 法を対角粘性項に採用しているため，線形拡散部に由来する陽的粘性 CFL は
等粘性流体では時間刻み決定に使用しない．
ただし，変粘性流体（気液二相流）ではクロス微分項 $\partial_y(\mu\partial_x v)$ が陽的処理されるため，
大密度比（$\mu_l/\mu_g \gg 1$）の場合には界面近傍で粘性類似の制約が実質的に復活しうる
（詳細は §~\ref{warn:cn_cross_derivative} の注意ボックス末尾の安定性解析を参照）．
\end{tcolorbox}

式~\eqref{eq:dt_sigma} の毛管波制約は，界面上の微小振幅波の分散関係
$\omega^2 = \sigma k^3/(\rho_l+\rho_g)$ にナイキスト波数 $k_{\max}=\pi/\Delta x$
を代入し，群速度 CFL 条件 $c_g\,\Delta t \le \Delta x$ を適用して得られる
（導出の全手順と $\rho_l\gg\rho_g$ での物理的解釈は付録~\ref{app:capillary_cfl} 参照）．
$\Delta t_\sigma \propto \Delta x^{3/2}$ の依存性は分散関係 $\omega\propto k^{3/2}$ に由来し，
格子を半分にすると約 2.8 倍の制約強化をもたらす．

\noindent\textbf{注意：}
表面張力が支配的な問題（$We \ll 1$，小液滴，毛管現象）では，
毛管波（capillary wave）の位相速度が大きくなり CFL 条件が厳しくなる．
高解像度計算では毛管波制約が支配的になることが多い．

\textbf{対策：}
\begin{itemize}
  \item 表面張力の半陰的処理（Implicit Surface Tension）による制約の緩和
  \item 問題によっては圧力 $p$ から $p + p_0$（静水圧 $p_0 = \rho g z$）を分離する
        動的圧力（kinematic pressure）分割で寄生流れを抑制しつつ $\Delta t$ を拡大
\end{itemize}

以上で，CCD エンジン（§~\ref{sec:CCD}），界面適合格子（§~\ref{sec:grid_gen}），
時間積分フレームワーク（§~\ref{sec:time_int}），
および CLS 移流・再初期化の空間離散化（本章）が揃った．
残る中核課題は，コロケート格子上での圧力・速度連成である．
次章（§~\ref{sec:collocate}）では，Balanced-Force 条件に基づく
離散レベルでの力の釣り合いの設計と，チェッカーボード不安定の抑制を論じる．
